\section{Experiments}\label{sec:experiments}
\subsection{Datasets and Few-Shot Splits}
We evaluate on MVTec AD~\cite{mvtec2019} and VisA~\cite{visa2022}; the strict external-transfer audit additionally uses MPDD. MVTec AD contains industrial object and texture categories with pixel-level defect masks. VisA contains a broader set of industrial categories and is used as the main full-scale corruption and calibration-shift benchmark. MPDD supplies six metal-part categories for an external MVTec-to-MPDD stress test; within-MPDD category certification is not claimed because its category count is too small for a three-way source split. For each class, support sets are sampled from normal training images with fixed seeds. Unless otherwise specified, $k\in\{1,2,4,8\}$ is used for clean accuracy-storage and strict nested studies, while the earlier conformal corruption tables use $k\in\{4,8\}$.

\subsection{Evaluation Protocols}
We organize experiments into six protocol groups. The first five cover clean accuracy/storage, PCA64/PCA128 sensitivity, corruption reliability, representative routing, and attainable operating levels. The sixth is the frozen strict nested protocol: $k\in\{1,2,4,8\}$, seeds 0--4, $\alpha\in\{0.05,0.10,0.20\}$, clean plus Gaussian noise, blur, brightness/contrast, and JPEG, at most 120 evaluation images per cell, $\rho=0.01$, PCA64, $\delta=0.05$, and at most 20 candidate thresholds. It evaluates MVTec and VisA within-dataset, MVTec-to-VisA, and MVTec-to-MPDD, with matched, clean-source, condition-agnostic, and mismatched negative-control routing. The pre-specified operational gate requires a nonzero category-certified threshold in at least 80\% of target cells, target FAR at most $\alpha+0.02$, nonzero power, at least 80\% no-harm, and a positive simultaneous power-gain lower bound. Zero thresholds and failed cells are retained.

\subsection{Naming Convention}
CRR denotes the overall framework (decoupled ranking plus reliability layer). The detector instantiations are named by their subspace configuration: CRR/PCA64 (default, also called CalibSubspaceHead in configuration files) and CRR/PCA128. Reliability views are named by their route: scalar/Vector/Shift-Aware/weighted Platt, temperature scaling, isotonic regression, histogram binning, LOIO conformal, and weighted conformal. Tables use these names consistently.

\subsection{Baselines and Ablations}
We compare against a controlled nearest-neighbor memory bank using DINOv2, a SubspaceAD-style PCA residual baseline, scalar Platt, Vector Platt, Shift-Aware Platt, weighted Platt, temperature scaling, isotonic regression, histogram binning, LOIO conformal, and weighted conformal. The controlled memory-bank implementation is not an official reproduction of PatchCore or AnomalyDINO; it holds the frozen DINOv2 feature pipeline and data split fixed to isolate the scoring and storage trade-off. In addition, we run the official AnomalyDINO code (ViT-S/14 at 448px, agnostic preprocessing) on both MVTec and VisA (official 1-cls split) for $k\in\{1,4,8\}$ with three seeds and report it as a separate accuracy reference. As the vision--language representative we include WinCLIP~\cite{winclip2023}, but only as reported numbers: no official implementation is released, and when we audited the leading community reimplementation under its own protocol (ViT-B-16+240, LAION-400M; zero-shot and 1-shot with its fixed support lists, three runs), it reproduced 71.7/77.4 mean image AUROC on MVTec and 66.0/70.0 on VisA against the reported 91.8/93.1 and 78.1/83.8---a roughly 20-point gap that the reimplementation's own authors also report and that we cannot attribute. This is precisely why our claim rules separate audited rows from reported rows, and WinCLIP therefore enters Table~\ref{tab:clean-efficiency} as a cited reference, not an audited baseline. Official SubspaceAD representative runs are included as a novelty guardrail. Their strong performance confirms that frozen DINOv2 subspace residuals are an existing competitive ranking mechanism; the residual is inherited, and CRR's contribution begins at the reliability of the alarms built on it.

\subsection{Metrics}
We separate ranking metrics from reliability metrics. AUROC and AP are computed from the raw PCA/subspace residual score $s(x)$. Pixel-level localization is measured by pixel AUROC and AU-PRO when masks are available. Reliability is measured operationally by empirical false-alarm rate at fixed $\alpha$, detection power, alarm precision, discrete uniformity of normal p-values, and risk--coverage curves; expected calibration error (ECE), Brier score, and negative log-likelihood (NLL) are reported as secondary metrics using the selected reliability probability or probability-like conformal view. Storage is measured as retained reference/model state per class. Robustness experiments include corruption shift and a label-aware PCA-residual FGSM surrogate; the latter is used only as a fragility diagnostic.

\subsection{Reproducibility and Claim Rules}
Experiments were run on a single NVIDIA RTX 5090 GPU with PyTorch 2.12.1 and DINOv2 patch features cached in fp32. The immutable handoff records the frozen commit, DINOv2 source/weight hashes, environment, dataset counts, per-image predictions, base-image identities, support manifests, corruption parameters, source partitions, candidate bounds, and SHA-256 checksums. Automated audits pass for 1,500 MVTec, 1,200 VisA, and 600 MPDD view cells and reject duplicate cells, missing support statistics, or support--test overlap. Reliability routes cannot claim ranking improvement unless the raw score changes; image-level and category-level certificates cannot substitute for one another.
